\chapter{\stateoftheartname}
\section{Narices Electrónicas (E-Noses)}

\subsection{Concepto}
Una nariz electrónica (e-Nose) es un dispositivo que se considera parte de la rama de investigación llamada biomímesis, ya que imita el sentido del olfato humano mediante el uso de sensores químicos para detectar y reconocer olores y compuestos volátiles en el aire. Estos dispositivos están diseñados para identificar patrones específicos de gases y vapores, permitiendo su aplicación en diversas áreas como la seguridad, la industria alimentaria, la medicina y el control ambiental.\\

Para que un dispositivo pueda considerarse una nariz electrónica debe cumplir varios requisitos fundamentales. En primer lugar, ha de ser capaz de detectar y diferenciar de forma simultánea múltiples compuestos volátiles; además, debe generar patrones de respuesta característicos que permitan su identificación mediante técnicas de análisis de datos.

Es imprescindible que ofrezca repetibilidad y estabilidad en sus mediciones a lo largo del tiempo y frente a distintas condiciones ambientales, y que pueda adaptarse a variaciones de temperatura, humedad y presión. Finalmente, la nariz electrónica no debe limitarse a proporcionar lecturas crudas de los sensores, debe entregar una salida procesada (mediante preprocesado y reconocimiento de patrones) que interprete y clasifique las señales para producir identificaciones o decisiones útiles y robustas.\\

Las principales ventajas de las narices electrónicas es que son rápidas, portatiles y económicas en comparación con metodos tradicionales para ciertas aplicaciones como la detección de sustancias, o técnicas cromatográficas para controles de calidad.

\subsection{Antecedentes y Evolución de la Nariz Electrónica}
La base tecnológica de la Nariz Electrónica (E-Nose) se encuentra en la investigación de los sensores de gas de semiconductores de óxido metálico (MOX), cuyo desarrollo se remonta a mediados del siglo XX.\@ Un hito temprano fue la descripción de las propiedades sensibles al gas del óxido de germanio en 1953. Los primeros sensores de gas se fabricaron utilizando óxido de zinc (\ce{ZnO}) y, a lo largo del tiempo, la tecnología evolucionó para incorporar materiales como el dióxido de estaño (\ce{SnO2}), el dióxido de titanio (\ce{TiO2}) y el óxido de tungsteno (\ce{WO3}). De estos, el dióxido de estaño granular es el material más ampliamente estudiado y empleado en la actualidad.

A pesar de que existen registros de sistemas de nariz electrónica reportados desde 1974, sus primeras aplicaciones prácticas se dieron en los laboratorios de control de calidad para el testeo de alimentos procesados. La formalización y comercialización de esta tecnología se materializó a principios de los años 90 con la aparición de los sistemas de reconocimiento de aroma en el mercado. Un sistema pionero fue el Modular Sensor System (MOSES I), marcando el inicio de la disponibilidad de estos instrumentos para una variedad de aplicaciones.

En la actualidad, esta tecnología ha evolucionado hasta convertirse en un instrumento capaz de reconocer y comparar olores complejos, un logro que la enmarca dentro del campo de la biomímesis al buscar imitar el olfato humano.

\subsection{Funcionamiento}
El funcionamiento de una Nariz Electrónica se basa en el principio de la biomímesis, buscando replicar el proceso del olfato humano mediante la combinación de hardware (la matriz de sensores) y software (el sistema de reconocimiento de patrones). El objetivo fundamental es la detección de moléculas volátiles en el aire y su posterior comparación con muestras de referencia calibradas.

El sistema completo se divide en tres partes esenciales que trabajan de forma secuencial: la toma de muestra (equipo de muestreo), la matriz de sensores y el sistema de tratamiento de datos.

\subsubsection{Ciclo de Medición y Detección}
El proceso se inicia con la preparación de la muestra, donde los Compuestos Orgánicos Volátiles (COV) son concentrados por calentamiento en la fase de vapor sobre la muestra (lo que se conoce como espacio de cabeza). Estos volátiles son introducidos en el sistema de detección a través de un ciclo de medición bien definido:
\begin{itemize}
    \item \textbf{Fase de Adsorción (Exposición):} Los sensores se exponen directamente a la atmósfera de la muestra contaminada (o espacio de cabeza), permitiendo que las moléculas de olor se adsorban sobre la superficie de los sensores. Durante esta fase, los sensores registran sus respuestas continuamente.
    \item \textbf{Fase de Desorción (Línea Base):} A continuación, los sensores son expuestos a aire limpio. Esto permite que las moléculas de olor se desorban, lo que restablece la línea base del sensor para la siguiente medición y es crucial para la posterior diferenciación de olores.
\end{itemize}

Los sensores de gas, principalmente los semiconductores de óxido metálico (MOX), miden las diferentes propiedades físico-químicas de los componentes del aroma. Esta interacción modifica las propiedades eléctricas del sensor (como su conductividad), convirtiendo el olor en una señal eléctrica mensurable.

\subsubsection{Generación de la Huella Olfativa}
La E-Nose requiere una matriz de múltiples sensores de gas para funcionar. Cada sensor posee una cierta especificidad para un grupo de compuestos químicos. Debido a que los olores suelen estar compuestos por más de una fuente, cada sensor de la matriz reacciona de forma única, generando un valor de conductividad individual.

El resultado de esta matriz de respuestas es una matriz de datos compleja, que constituye la ``huella digital'' o ``huella olfativa'' del olor percibido. Esta huella es esencial para la identificación y clasificación de los olores, ya que cada combinación de respuestas de los sensores corresponde a un aroma específico.

\subsubsection{Procesamiento y Reconocimiento de Patrones}
La etapa final y definitoria del sistema es el procesamiento de datos. El ordenador se encarga de procesar la matriz de datos mediante técnicas quimiométricas (estadísticas multivariadas y algoritmos de aprendizaje automático).

Métodos avanzados como el Análisis de Componentes Principales (PCA), el Análisis de Función de Discriminantes (DFA) o las Redes Neuronales Artificiales (ANN) son esenciales para la extracción de patrones significativos y la clasificación del olor. Es precisamente este procesamiento algorítmico lo que cumple el criterio de ``producir un resultado que es diferente a solo las respuestas de los sensores de gas'', ya que permite a la E-Nose clasificar los diferentes elementos en función de su similitud aromática e Identificar si el olor se ajusta a un patrón conocido.

De esta forma, la nariz electrónica es capaz de procesar la fracción volátil de forma global, al igual que la nariz humana, y proporciona una interpretación sobre la naturaleza del aroma.

\subsection{Tipos de sensores}
Para la detección de olores, las narices electrónicas emplean diversos tipos de sensores, cada uno con sus propias características y aplicaciones. Los tipos de sensores más comunes son:
\begin{itemize}
    \item Semiconductores de óxido metálico (MOS o MOX).
    \item Transistores de campo eléctrico con semiconductores de óxido metálico (MOSFET).
    \item Conductores de polímero orgánico (CP).
    \item Cristales piezoeléctricos (BAW).
    \item Detectores de fotoionización, biosensores y nanotecnología.
\end{itemize}

La vida útil de los sensores MOSFET y CP, típicamente de 1 a 3 años, se ve limitada por su alta dependencia de la temperatura y la humedad, a pesar de ofrecer una sensibilidad aceptable.

Los cristales piezoeléctricos (BAW), que funcionan como microbalanzas de alta precisión, ofrecen una gran sensibilidad a la masa de las moléculas adsorbidas, pero su rendimiento se ve comprometido por la humedad y la temperatura. Por otro lado, los detectores de fotoionización proporcionan una respuesta casi instantánea y son muy sensibles a compuestos orgánicos volátiles, aunque carecen de selectividad. En un intento por maximizar la especificidad, los biosensores integran elementos biológicos (como enzimas o anticuerpos), pero su fragilidad y corta vida útil limitan su aplicación práctica fuera del laboratorio. Finalmente, la nanotecnología, mediante el uso de materiales como nanotubos de carbono o grafeno, promete una sensibilidad sin precedentes y un consumo energético mínimo, si bien la reproducibilidad en su fabricación a gran escala sigue siendo un desafío.

La clase de sensores mas utilizada son los sensores semiconductores de óxido metálico (MOX). De hecho, más del 35\% de los sistemas de nariz electrónica reportados desde 1974 se basan en la tecnología MOX.\@ Los sensores MOX tienen selectividad y reproducibilidad pobres, pero son poco dependientes de la humedad, lo que les da una vida útil de entre 3 y 5 años. Funcionan cambiando la conductividad eléctrica del semiconductor (como \ce{SnO2} o \ce{TiO2}) cuando interacciona con los compuestos volátiles. 

El principio radica en un proceso de quimisorción, que es la formación de enlaces químicos entre las moléculas de gas y el semiconductor.La quimisorción altera la concentración de los portadores de corriente en la superficie del semiconductor, lo que resulta en un cambio en su conductividad eléctrica. Finalmente, cuando el sensor se encuentra en una atmósfera de impurezas gaseosas reductoras, la barrera de energía disminuye. Esto, desde un punto de vista práctico, significa que la resistencia (R) disminuye cuando el sensor se expone a aire contaminado.

La resistencia del sensor puede expresarse mediante la fórmula:

\begin{equation}
    R = R_Z \cdot \exp\left(\frac{W_G}{k_B T}\right)
    \label{eq:resistencia_sensor}
\end{equation}

Donde $R_Z$ representa la resistencia interna al grano ($\Omega$), $W_G$ es la barrera energética (eV), $k_B$ es la constante de Boltzmann ($eV \cdot K^{-1}$), y T es la temperatura absoluta (K).`\\

Los sensores MOX se clasifican según la técnica utilizada para producir la capa sensible al gas. Por un lado están los sensores de película delgada (Thin-Film), cuyo espesor de la capa oscila entre 5nm y 2$\mu$m. Por otra parte están los sensores de película gruesa (Thick-Film), cuyo espesor de la capa es de 10 a 300 $\mu$m. En comparación con los de película delgada, son menos susceptibles al envenenamiento, son más duraderos y exhiben una mayor sensibilidad.

\subsubsection{Tipos de materiales MOX}
Los sensores MOX originales estaban hechos de óxido de zinc (\ce{ZnO}) y fueron seguidos por el uso de otros materiales. Los óxidos de metal que se han utilizado o propuesto se muestran en la Tabla~\ref{tab:materiales_mox}.

\begin{table}[htbp]
\centering
\begin{tabularx}{\textwidth}{lX}
\toprule
\textbf{Material} & \textbf{Aplicaciones y Notas Clave} \\
\midrule
Dióxido de Estaño (\ce{SnO2}) & El más ampliamente estudiado y utilizado. Se usa en sensores de película porosa y gruesa con una alta relación superficie-volumen. \\
\midrule 
Dióxido de Titanio (\ce{TiO2}) & Compuesto generalizado, no tóxico, barato y respetuoso con el medio ambiente. Junto con \ce{SnO2}, puede dar lugar a sensores de alta selectividad. \\
\midrule 
Óxido de Zinc (\ce{ZnO}) & Utilizado para la detección de gases, monitoreo ambiental, procesos industriales y análisis de respiración. Se utiliza para monitorear la calidad del aire interior. \\
\midrule 
Óxido de Tungsteno (\ce{WO3}) & Utilizado en detección de gas, monitoreo ambiental, procesos industriales y en el desarrollo de ventanas inteligentes. \\
\midrule 
Óxido de Indio (\ce{In2O3}) & Se utiliza para monitorear y regular contaminantes como óxidos de nitrógeno (\ce{NO_x}), óxidos de azufre (\ce{SO_x}) y Compuestos Orgánicos Volátiles (COV). \\
\midrule 
Óxido Cúprico (\ce{CuO}) & Ha sido utilizado para diversas aplicaciones como baterías, biosensores, células solares y sensores de gas. \\
\midrule 
Óxido Cuproso (\ce{Cu2O}) & Abordado menos en aplicaciones de detección de gas en comparación con \ce{CuO}. \\
\bottomrule
\end{tabularx}
\caption{Materiales comunes en sensores MOX y sus aplicaciones.}\label{tab:materiales_mox}
\end{table}

\subsubsection{Desafíos y Mantenimiento}

\subsection{Procesamiento de Datos (Quimiometría)}

\subsection{Aplicaciones y Desafíos de E-Noses}

\section{Agentes caninos}
